\documentclass[10pt,a4paper]{amsart}
\usepackage[margin=19mm]{geometry}
\usepackage{amsmath,amssymb,mathtools}
\usepackage[scr=rsfs,cal=euler]{mathalfa}
\usepackage{xfrac}
\usepackage[unicode,hidelinks,bookmarks=false]{hyperref}
\newcommand{\ie}{i.e.\ }
\newcommand{\cf}{cf.\ }
\newcommand{\resp}{respectively\ }
\newcommand{\Subsection}{\subsection}
\providecommand{\nobreakdash}{\nobreak\hbox{-}}
\setlength{\emergencystretch}{2em}
\multlinegap0pt
\usepackage{amssymb}
\usepackage{xcolor}
\newtheorem{thm}{Theorem}
\newtheorem{lem}{Lemma}
\newtheorem{cor}{Corollary}
\newcommand{\NN}{\mathbb{N}}
\newcommand{\PP}{\mathbb{P}}
\newcommand{\Pa}{\mathcal{P}}
\newcommand{\RR}{\mathbb{R}}
\newcommand{\dtv}{\mathrm{d}_{\mathrm{TV}}}
\title{A Kubilius model for sieve-theoretic sequences}
\author{Ofir Gorodetsky}
\date{Source: arXiv:2608.14190v1 (CC BY 4.0)}
\begin{document}
\maketitle

\noindent\emph{Source and licence.} A Kubilius model for sieve-theoretic sequences. Version arXiv:2608.14190v1 (CC BY 4.0), distributed by arXiv under the Creative Commons Attribution 4.0 International licence, http://creativecommons.org/licenses/by/4.0/, as recorded in the arXiv licence metadata for this exact version and shown on the abstract page https://arxiv.org/abs/2608.14190v1. Full licence notice: \texttt{licenses/arxiv-2608.14190-CC-BY-4.0.txt}.\par\smallskip

\noindent\textit{Editorial scope.} Counted unit: the article's cor cor:rankin (source0.tex lines 402--408). The article's complete original proof of it is delivered with it (source0.tex lines 409--418, one \texttt{proof} environment of 10 lines). No mathematics is written, repaired or re-derived: the statement and the proof are the article's own lines, in the article's own notation, and no retained line is substituted or re-flowed.  Retained uncounted as the source-local context the proof needs: lines 166--187; lines 360--369.  Nothing was omitted from any retained span, so the package carries no omission record.  No retained line contains a citation, so the wrapper carries no bibliography and no bibliography entry.  No retained line contains a figure environment, an image-inclusion command or drawing code, so no image is delivered and the package declares no asset dependency.  Editorial formatting only, in addition: an AMS \texttt{amsart} preamble that restates the article's own declarations and macros used by the retained lines, namely the environments \texttt{thm}, \texttt{lem}, \texttt{cor} on separate counters, together with the standard packages those lines need.  The retained environments print in this document as Theorem 1, Lemma 1, Corollary 1 in order of appearance, whereas the article prints the same items with the numbering of its own full document.  The complete native source of the article as distributed in the arXiv e-print for this exact version is bundled unchanged as \texttt{sources/207643/source0.tex}.
\begin{thm}\label{thm:KM}
	Consider the following input:
	\begin{itemize}
		\item Let $A$ be a random variable  taking values in $[1,x]\cap \NN$. 
		\item Let $y\in [2,x]$.
			\item Let $\Pa$ be a subset of the primes up to $y$.
		\item Let $g\colon \NN\to [0,1]$ be a multiplicative function with $g(d)=0$ for all $d\not\in \langle \Pa\rangle$. Suppose that $g(p^{k+1})\le g(p^k)$	for all $p\in \Pa$ and $k\in \NN_{\ge 0}$, and that $\lim_{k\to \infty}g(p^k)=0$ for all $p\in \Pa$. Suppose that there exist $\kappa\ge 0$ and $B>0$ such that the following holds: for every $d\in \langle\Pa\rangle$ such that $g(p^{\nu_p(d)+1})<g(p^{\nu_p(d)})$ for all $p \in \Pa$, we have
		\begin{equation}\label{eq:iwaniecd}
			\prod_{a\le p \le b} \left( 1-\frac{g(pd)}{g(d)}\right)^{-1} \le \left( \frac{\log b}{\log a}\right)^{\kappa}\exp\left( \frac{B}{\log a}\right)
		\end{equation}
	for all $2\le a \le b\le y$.
		\item Let $(G_p)_{p\in \Pa}$ be independent random variables taking values in $\NN_{\ge 0}$ with $\PP(G_p=k)=g(p^k)-g(p^{k+1})$ for all $k\in \NN_{\ge 0}$.
	\end{itemize}
	Then, 	for any parameters $S>T>0$ with $S/T\ge y$, we have
	\[	\dtv( (\nu_p(A))_{p\in\Pa}, (G_p)_{p\in \Pa}) \ll_{\kappa,B} D_1+D_2\]
	for
	\begin{align*}
		D_1 &=\sum_{\substack{ d>T\\ g(d)\neq 0}} g(d)\prod_{p\in \Pa}(1-g(p^{\nu_p(d)+1})/g(p^{\nu_p(d)}))+\sum_{\substack{ d \le T\\ g(d)\neq 0}} g(d) \prod_{p\in \Pa}(1-g(p^{\nu_p(d)+1})/g(p^{\nu_p(d)})) f(u_d),\\
		D_2 &=  \sum_{\substack{d\le S\\ d \in \langle \Pa\rangle }} |\PP(d\mid A)-g(d)|(2^{\omega(d)}+\mathbf{1}_{g(d)\neq 0}\prod_{p\in \Pa}(1+g(p^{\nu_p(d)+1})/g(p^{\nu_p(d)}))),
	\end{align*}
	where $u_d$ is given by $y^{u_d} := S/d$, and $f(t):=e^{-t\log t+t\log_3 \max\{t,100\}+C_{\kappa,B}t}$ for a sufficiently large constant  $C_{\kappa,B}$. Here $\log_3 t=\log \log \log t$.
\end{thm}

\begin{lem}\label{lem:rankin}
Let $h\colon \NN\to \RR_{\ge 0}$ be a nonnegative multiplicative function. Suppose that for some $\kappa,B> 0$ we have
	\begin{equation}\label{eq:mertensk}
		\bigg|\sum_{p\le t} \frac{h(p)\log p}{p}-\kappa \log t\bigg|\le B
	\end{equation}
	for all $t \in [2,y]$. Suppose $x\ge y \ge 2$. If $h$ is supported on squarefrees then
	\[ \sum_{\substack{d>x\\ P(d)\le y}} \frac{h(d)}{d} \ll \prod_{p\le y}\left(1+\frac{h(p)}{p}\right) e^{O_{\kappa,B}(u)}  (u\log (u+1))^{-u}. \]
	If for all $p\le y$ and $i\ge 2$ we have $h(p^i)\le (i+1)^B$, then
	\[ \sum_{\substack{d>x\\ P(d)\le y}} \frac{h(d)}{d} \ll  \prod_{p\le y}\left(\sum_{i=0}^{\infty}\frac{h(p^i)}{p^i}\right)\cdot  \begin{cases}  (u\log (u+1))^{-u}e^{O_{\kappa,B}(u)} &\text{if }x\ge y \ge \log x,\\ x^{-1} e^{O_{\kappa,B}\big(\log \big(1+\tfrac{\log x}{y}\big)\tfrac{y}{\log y}\big)} &\text{if }\log x \ge y \ge 2.\end{cases}\]
\end{lem}

\begin{cor}\label{cor:rankin}
	Let $x\ge y\ge 2$. Under the same notation and assumptions of Lemma \ref{lem:rankin}, we have
	\[ \sum_{\substack{d\le x\\ P(d)\le y}} \frac{h(d)}{d}f\big( u+1 - \frac{\log d}{\log y}\big)\ll \prod_{p\le y}\left(1+\frac{h(p)}{p}\right) u^{-u}e^{O_{\kappa,B}(u\log_3 \max\{u,100\})}\]
for $x\ge y \ge 2$	if $h$ is supported on squarefrees, and 
	\[ \sum_{\substack{d\le x\\ P(d)\le y}} \frac{h(d)}{d}f\big( u+1 - \frac{\log d}{\log y}\big)\ll\prod_{p\le y}\left(\sum_{i=0}^{\infty}\frac{h(p^i)}{p^i}\right) \cdot \begin{cases} u^{-u}e^{O_{\kappa,B}(u\log_3 \max\{u,100\})} &\text{if }x\ge y \ge \log x, \\ x^{-1+O_{\kappa,B}(\log_3 \max\{x,100\} / \log \log x)} &\text{if }\log x \ge y \ge 2\end{cases}\]
	if $h(p^i)\le (i+1)^B$ for $i\ge 2$. Here $f$ is as in Theorem \ref{thm:KM}.
\end{cor}

\begin{proof}
We may discard the $d=1$ term. By considering $d\in (y^{\ell},y^{\ell+1}]$ for each $0\le \ell \le u$, we have
\begin{align}
\notag\sum_{\substack{1<d\le x\\ P(d)\le y}} \frac{h(d)}{d}f\big( u+1 - \frac{\log d}{\log y}\big) &\le \sum_{0\le \ell \le u} \big(\max_{d \in (y^{\ell},y^{\ell+1}]} f\big( u+1 - \frac{\log d}{\log y}\big)\big) \sum_{\substack{d>y^{\ell}\\ P(d)\le y}} \frac{h(d)}{d}\\
\label{eq:hmax}&\le (u+1) \max_{0\le \ell \le u} \big(\max_{d \in (y^{\ell},y^{\ell+1}]} f\big( u+1 - \frac{\log d}{\log y}\big)\big) \sum_{\substack{d>y^{\ell}\\ P(d)\le y}} \frac{h(d)}{d}. 
\end{align}
The $(u+1)$-factor is harmless. From properties of $f$, and Lemma~\ref{lem:rankin} with $x=y^{\ell}$, we can estimate the maximum over $d$ as well as the $d$-tail in \eqref{eq:hmax}. If $h$ is supported on squarefrees, we obtain that the right-hand side of \eqref{eq:hmax} is at most
\[  \prod_{p\le y}(\sum_{i=0}^{\infty}\frac{h(p^i)}{p^i}) e^{O_{\kappa,B}(u\log_3 \max\{u,100\})} \max_{0\le \ell \le u}((\ell+1)\log(\ell+2))^{-\ell}(u-\ell+1)^{-(u-\ell)}. \]
The maximum is of size $\ll u^{-u} e^{O(u)}$, concluding the proof in this case. If $h(p^i)\le (1+i)^B$ is assumed for $i\ge 2$, the same argument works when $y\ge \log x$. If $y<\log x$ one needs to consider $u\ge \ell > y/\log y$ separately from $\ell \le y/\log y$ because of the two ranges in Lemma~\ref{lem:rankin}. The maximum over $\ell\le y/\log y$ contributes at most $u^{-u}e^{O_{\kappa,B}(u\log_3 \max\{u,100\})}$, which is acceptable. The maximum over $u\ge \ell >y/\log y$ contributes $\ll x^{-1+O_{\kappa,B}(\log_3 \max\{x,100\} / \log \log x)}$ by a calculus computation that is left to the reader.
\end{proof}

\end{document}
